\section*{Hoofdstuk \ref{ch:b2:fragment-orbitals} --- Fragmentorbitalen en meeratomige moleculen}

\begin{solution}{exo:b2:fragment-orbitals:1}
$h_1 + h_2$ en $h_1 - h_2$. Het vlak $xz$, dat de hoek middendoor deelt en
loodrecht op het molecuul staat, verwisselt de twee waterstofatomen: $h_1 - h_2$ verandert
van teken, $h_1 + h_2$ niet.
\end{solution}

\begin{solution}{exo:b2:fragment-orbitals:2}
Zes valentieorbitalen (de 2s en de drie 2p van zuurstof, de twee 1s van waterstof), acht valentie-elektronen,
vier gevulde orbitalen.
\end{solution}

\begin{solution}{exo:b2:fragment-orbitals:3}
De gevulde valentieniveaus zijn één $a_1$-orbitaal en een verzameling van drie ontaarde
$t_2$-orbitalen: twee energieën, twee banden. De $t_2$-band verwijdert elektronen uit drie
orbitalen met zes elektronen, tegenover één orbitaal en twee elektronen voor $a_1$.
\end{solution}

\begin{solution}{exo:b2:fragment-orbitals:4}
De $\pi$-binding vereist dat de twee p-orbitalen van de koolstofatomen parallel zijn, wat
beide \ce{CH2}-vlakken in één gemeenschappelijk vlak legt.
\end{solution}

\begin{solution}{exo:b2:fragment-orbitals:5}
Ammoniak verliest een elektron uit zijn vrij elektronenpaar, een overwegend \omterm{def:b2:diatomic-mos:bonding}{niet-bindende orbitaal} op
stikstof; methaan uit een bindende $t_2$-orbitaal. Een niet-bindend elektron ligt hoger,
dichter bij het atomaire niveau, dan een bindend: het wordt gemakkelijker verwijderd.
\end{solution}

\begin{solution}{exo:b2:fragment-orbitals:6}
Lineair water heeft twee ontaarde, gevulde, niet-bindende p-orbitalen; buigen verlaagt er één
van en verhoogt een bindend niveau licht: met acht elektronen een netto winst. \ce{BeH2}
heeft er maar vier, die de twee \omterm{def:b2:diatomic-mos:bonding}{bindende orbitalen} vullen; de orbitaal die door het buigen
zou dalen, is leeg, terwijl de bindende niveaus overlap verliezen: de lineaire vorm is de beste.
\end{solution}

\begin{solution}{exo:b2:fragment-orbitals:7}
Elke C--H- of C--C-binding op een koolstofatoom naast het kationische centrum kan zich richten naar de
lege p-orbitaal en levert ongeveer $2\beta^2/\Delta$ stabilisatie. \ce{CH3+} heeft geen zulke
buur; het ethyl-, isopropyl- en tert-butylkation hebben één, twee en drie methylgroepen
naast het centrum, die elk een goed gerichte binding bieden: de stabiliteit groeit in die
volgorde.
\end{solution}

\begin{solution}{exo:b2:fragment-orbitals:8}
Voor twee gelijke niveaus is de splitsing $2|\beta|$ en winnen de twee $\pi$-elektronen $2|\beta|$,
evenredig met $\cos\theta$. Ze daalt tot de helft bij $\theta = 60^\circ$ en tot nul bij
$90^\circ$.
\end{solution}

\begin{solution}{exo:b2:fragment-orbitals:9}
Boraan heeft zes valentie-elektronen: ze vullen de drie \omterm{def:b2:diatomic-mos:bonding}{bindende orbitalen}, en de p-orbitaal
van boor loodrecht op het vlak blijft leeg; piramidaliseren levert niets op.
Ammoniak heeft er twee meer: ze bezetten de orbitaal die in de vlakke vorm niet-bindend is, en,
zoals in water, laat piramidaliseren haar mengen met de 2s-orbitaal en
de waterstofcombinatie, waardoor ze daalt.
\end{solution}

\begin{solution}{exo:b2:fragment-orbitals:10}
De seculiere vergelijkingen voor $c_1 = c_2 = c_3$ geven $\alpha + 2\beta$; de twee combinaties
die er orthogonaal op staan, geven $\alpha - \beta$ (ontaard). Twee elektronen in $\alpha + 2\beta$: energie
$2\alpha + 4\beta$, lager dan $\ce{H2} + \ce{H+}$ ($2\alpha + 2\beta$), want $\beta < 0$.
\end{solution}

\begin{solution}{exo:b2:fragment-orbitals:11}
$\pi_1$: alle drie coëfficiënten van hetzelfde teken, geen knoop (bindend); $\pi_2$: tegengestelde tekens
op de eindstandige koolstofatomen en een knoop in het midden (niet-bindend); $\pi_3$: afwisselende
tekens, twee knopen (antibindend). Het anion heeft vier $\pi$-elektronen: $\pi_2$ is zijn
hoogste bezette niveau, met alleen dichtheid op de eindstandige koolstofatomen.
\end{solution}

\begin{solution}{exo:b2:fragment-orbitals:12}
In elk van de twee vlakken die de as bevatten, geven drie parallelle p-orbitalen (O, C, O)
een bindende, een niet-bindende (alleen op de zuurstofatomen) en een \omterm{def:b2:diatomic-mos:bonding}{antibindende orbitaal}: zes
$\pi$-orbitalen in totaal. De twee bindende en de twee niet-bindende zijn gevuld: acht
$\pi$-elektronen (twee $\pi$-bindingen en twee vrije elektronenparen van zuurstof in het lewisbeeld).
\end{solution}

\begin{solution}{pb:b2:fragment-orbitals:1}
\textbf{1.} Acht: zes van zuurstof, één van elk waterstofatoom.
\textbf{2.} $h_1 + h_2$ en $h_1 - h_2$.
\textbf{3.} 2s en $2\mathrm p_z$.
\textbf{4.} $2\mathrm p_y$.
\textbf{5.} $2\mathrm p_x$, loodrecht op het molecuulvlak.
\textbf{6.} Zes; vier gevuld.
\textbf{7.} $h_1 - h_2$, die, zoals $2\mathrm p_z$, van teken verandert bij het vlak
loodrecht op de as door de zuurstof.
\textbf{8.} $h_1 + h_2$.
\textbf{9.} $2\mathrm p_x$ en $2\mathrm p_y$, loodrecht op de as: een ontaard
niet-bindend paar.
\textbf{10.} (2s, $h_1 + h_2$) bindend$^2$, ($2\mathrm p_z$, $h_1 - h_2$) bindend$^2$, daarna
$2\mathrm p_x^2\,2\mathrm p_y^2$.
\textbf{11.} Nee: het ontaarde paar is vol, elk elektron is gepaard.
\textbf{12.} Op de energie van een 2p-orbitaal van zuurstof: het is niet-bindend.
\textbf{13.} $2\mathrm p_y$ (met de as van het lineaire molecuul langs $z$ en buigen in $yz$).
\textbf{14.} Ze mengt met 2s en $h_1 + h_2$ en daalt: het niveau $3a_1$.
\textbf{15.} Het bindende niveau ($2\mathrm p_z$, $h_1 - h_2$), dat wat overlap verliest wanneer de
waterstofatomen de as verlaten: het wordt $1b_2$.
\textbf{16.} Het gevulde niveau dat daalt, wint meer dan het andere verliest: met acht
elektronen is de gebogen vorm de stabielste.
\textbf{17.} $104.48^\circ$, tussen $90^\circ$ (bindingen uit alleen p-orbitalen) en $109.5^\circ$:
de 2s-orbitaal neemt deel aan de binding.
\textbf{18.} $2\mathrm p_x$, die de niet-bindende $1b_1$ blijft.
\textbf{19.} Vier.
\textbf{20.} Uit $1b_1$, de niet-bindende $2\mathrm p_x$-orbitaal van zuurstof.
\textbf{21.} Een niet-bindend elektron verwijderen verandert de bindingen nauwelijks: het ion heeft de
geometrie van het molecuul en er wordt weinig vibratie-energie aangeslagen. Een
bindend elektron verwijderen verandert de bindingslengten en -hoeken, en de band spreidt zich over
veel vibratieniveaus.
\textbf{22.} \qty{12.62}{eV} tegenover \qty{13.62}{eV}: zuurstof draagt in water een negatieve
partiële lading die van de waterstofatomen getrokken is, wat zijn 2p-niveaus verhoogt.
\textbf{23.} Het zijn geen twee gelijke niveaus: hun twee combinaties zijn de orbitalen $3a_1$ en
$1b_1$, met verschillende energieën. Beide beelden beschrijven dezelfde totale
elektronendichtheid; alleen de gedelokaliseerde orbitalen geven de ionisatie-energieën.
\textbf{24.} Water heeft \textbf{vier} valentie-ionisatiebanden, de laagste bij
$\boldsymbol{\approx \qty{12.6}{eV}}$.
\end{solution}
